\needspace{6\baselineskip}\noindent \textbf{C.1. Ba nguyên tắc an toàn bắt buộc}

\begin{enumerate}[leftmargin=1.6em,itemsep=0.15em,topsep=0.3em]
  \item \textbf{Không cấp 5 V cho module LoRa.} Chỉ 3,3 V.
  \item \textbf{Không cấp nguồn cho module LoRa khi chưa gắn ăng-ten.} Việc này có thể phá hỏng tầng khuếch đại công suất.
  \item \textbf{Đo điện áp đầu ra của cả hai mạch hạ áp trước khi nối vào bất kỳ board nào.} XL4015 phải cho 6,00 $\pm$ 0,1 V; LM2596 phải cho 5,00 $\pm$ 0,1 V.
\end{enumerate}
\needspace{6\baselineskip}\noindent \textbf{C.2. Trình tự kiểm tra sau khi đấu nối}

Kê xe lên cao để bánh không chạm đất, nạp firmware, mở cổng nối tiếp. Trong lúc firmware đo độ trôi con quay lúc khởi động, \textbf{giữ xe đứng yên hoàn toàn}. Sau đó chạy lần lượt: lệnh tự kiểm tra tổng thể $\rightarrow$ kiểm tra công tắc va chạm $\rightarrow$ kiểm tra encoder $\rightarrow$ kiểm tra ADC của MQ-3 $\rightarrow$ chạy từng bên động cơ $\rightarrow$ thử đường truyền LoRa.

Chỉ khi mọi khối đã hoạt động đúng mới đặt xe xuống sàn để hiệu chuẩn theo thứ tự: quãng đường encoder $\rightarrow$ độ trôi con quay $\rightarrow$ lệnh đi thẳng 100 cm $\rightarrow$ lệnh quay 360$^\circ$ $\rightarrow$ đo độ lệch khi đi thẳng $\rightarrow$ hiệu chuẩn MQ-3 $\rightarrow$ chạy thuật toán tự hành.

\needspace{6\baselineskip}\noindent \textbf{C.3. Quy trình một buổi thí nghiệm}

\begin{enumerate}[leftmargin=1.6em,itemsep=0.15em,topsep=0.3em]
  \item Bật nguồn và \textbf{để cảm biến MQ-3 sấy nóng ít nhất 5 phút} trước khi tin vào số liệu.
  \item Đặt xe ở góc xuất phát, giữ yên trong lúc firmware đo giá trị nền và độ trôi con quay.
  \item Đặt nguồn phát (đĩa cồn) tại vị trí đã chọn; với thí nghiệm có gió, bật quạt và để dòng khí ổn định vài chục giây.
  \item Chọn thuật toán bằng nút trên board, rồi nhấn để bắt đầu.
  \item Theo dõi từ trạm mặt đất; chương trình trên máy tính ghi log thành tệp CSV kèm nhãn thuật toán, môi trường và số thứ tự lần chạy.
  \item Sau khi robot kết luận, đo khoảng cách thật từ đầu dò tới nguồn và ghi vào tệp dữ liệu đi kèm.
  \item Trước khi chạy lần tiếp theo, \textbf{tắt quạt, thông gió phòng và chờ giá trị nền trở về mức cũ} --- nếu không, giá trị nền của lần chạy sau sẽ bị nâng lên bởi lượng khí còn lại của lần chạy trước.
\end{enumerate}
